% -----------------------------------------------
% Resumo do TP1 (DCC831 IA Generativa para Música, UFMG 2026/2)
% Template ISMIR 2026. Limite: 2 páginas de conteúdo (graduação);
% referências e AI Usage Statement não contam.
%
% Convenção do esqueleto: \nota{...} (azul) marca o que falta escrever.
% Apague cada \nota quando o parágrafo estiver pronto. Antes de entregar,
% `grep -n nota resumo.tex` deve voltar vazio (fora a definição).
% -----------------------------------------------

\documentclass{article}
\usepackage[T1]{fontenc}
\usepackage[utf8]{inputenc}
\usepackage[brazil]{babel}
\usepackage{ismir} % sem "submission": o TP não é anônimo
\usepackage{amsmath,amssymb,cite,url}
\usepackage{graphicx}
\usepackage{color}
\usepackage{booktabs}
\usepackage{forest}

% O ismir.sty imprime no rodapé da 1ª página uma nota de copyright que atribui
% o texto aos anais do ISMIR 2026 (Abu Dhabi). O TP não foi publicado lá, então
% a nota é suprimida. Apague a linha abaixo para restaurá-la.
\renewcommand{\permission}{}

\newcommand{\nota}[1]{\textcolor{blue}{[#1]}}
\newcommand{\nt}[1]{\textsc{#1}}      % não-terminal
\newcommand{\tm}[1]{\texttt{#1}}      % terminal

% ------
\title{Uma Gramática Probabilística para Backing Tracks de Samba}

\oneauthor
  {Vinícius Pinho}
  {Departamento de Ciência da Computação, UFMG\\\texttt{viniciuspgalvao@ufmg.br}}

\def\authorname{V. Pinho}

\sloppy

\begin{document}

\maketitle

% =====================================================================
\begin{abstract}
Este trabalho gera \emph{backing tracks} de percussão de samba para quem estuda cavaquinho,
com uma gramática livre de contexto probabilística cujos terminais são ciclos de dois
compassos com todos os instrumentos. A restrição estilística é o samba em 2/4, com o ciclo
ancorado no telecoteco do tamborim e blocos de 16 compassos terminados por virada ou breque. A direção artística foi uma levada estável que varia de
forma contida e marca a forma; os hiperparâmetros controlam quanto tamborim e pandeiro
variam, quanto a cuíca aparece e quanto a música dura. Cada execução gera um MIDI
reprodutível pela seed e um áudio feito por um sampler próprio, com recortes do dataset BRID. A linguagem gerada é regular, e a gramática se
justifica pela legibilidade da hierarquia musical.
\end{abstract}

% =====================================================================
\section{Introdução}\label{sec:introducao}

Quem está aprendendo cavaquinho provavelmente já estudou sobre \emph{backing tracks}: faixas
com a percussão de uma roda de samba sobre as quais o instrumentista toca. A direção artística deste trabalho foi uma faixa com levada estável (a âncora), que
varie o suficiente para não soar como loop e que \emph{sinalize a forma}, para que quem toca
saiba onde está. A ideia é gerar faixas assim com variações sorteadas, para que o aluno
treine a integração com a percussão, o ritmo num andamento fixo e a improvisação, seguindo o
fluxo que a música anuncia.

A restrição estilística é o samba em 2/4, o compasso do gênero. O ciclo rítmico tem 2
compassos e se ancora no telecoteco do tamborim, a ``clave'' do samba~\cite{sandroni2001}.
Os ciclos formam frases de 8 compassos, agrupadas em blocos de 16, e o fim de cada bloco é
marcado por uma virada ou um breque. Os instrumentos são surdos de primeira e de segunda,
pandeiro, tamborim e cuíca. O método é uma gramática livre de contexto probabilística em que a hierarquia da forma
(música $\to$ bloco $\to$ frase $\to$ ciclo) está nas próprias regras, e cujas
probabilidades são os hiperparâmetros.

% =====================================================================
\section{Trabalhos Relacionados}\label{sec:relacionados}

O Brazilian Rhythmic Instruments Dataset (BRID)~\cite{maia2018} reúne gravações em estúdio
de percussão brasileira (274 faixas solo e 93 mixagens) para análise de ritmo. Aqui ele é a
fonte do timbre: as faixas solo de samba de um mesmo
instrumentista fornecem os recortes de surdo, pandeiro, tamborim e cuíca com que o nosso
sampler transforma o MIDI gerado em áudio (Seção~\ref{sec:audio}).

% =====================================================================
\section{Método}\label{sec:metodo}

\subsection{Representação}\label{sec:representacao}

O tempo é discretizado em semicolcheias: um compasso de 2/4 tem 8 passos. O terminal da
gramática é um \emph{ciclo} de 2 compassos (16 passos), a unidade em que o telecoteco do
tamborim se completa; um terminal de um compasso partiria a clave ao meio. Cada ciclo atribui
a cada uma das 8 pistas (instrumento e articulação) um valor por passo: 0 (silêncio),
1 (batida) ou 2 (acento), como na Figura~\ref{fig:grade}.

\begin{figure}[t]
  \centering
  \footnotesize
  \begin{tabular}{@{}lll@{}}
    \toprule
    pista            & \tm{base}                        & \tm{virada} \\
    \midrule
    surdo 1          & \texttt{....X...\,....X...} & \texttt{....X...\,X..X..X.} \\
    surdo 2          & \texttt{x.......\,x.......} & \texttt{x.......\,X..X..X.} \\
    pand. grave      & \texttt{x...x...\,x...x...} & \texttt{x...x...\,........} \\
    pand. platinela  & \texttt{xxxXxxxX\,xxxXxxxX} & \texttt{xxxXxxxX\,........} \\
    pand. tapa       & \texttt{.......x\,.......x} & \texttt{.......x\,X..X..X.} \\
    tamborim         & \texttt{x.x.xx.x\,.x.xx.x.} & \texttt{x.x.xx.x\,X..X..X.} \\
    cuíca grave      & \texttt{........\,........} & \texttt{........\,........} \\
    cuíca aguda      & \texttt{........\,........} & \texttt{........\,X..X..X.} \\
    \bottomrule
  \end{tabular}
  \caption{Dois terminais na grade de 16 semicolcheias (um ciclo = 2 compassos de 2/4).
  \texttt{.}~silêncio, \texttt{x}~batida, \texttt{X}~acento. Na \tm{virada}, o 2º compasso
  troca a levada por uma figura 3-3-2 em uníssono: acentos nos passos 0, 3 e 6, que dividem
  as 8 semicolcheias em grupos de 3, 3 e 2.}
  \label{fig:grade}
\end{figure}

\subsection{Gramática}\label{sec:gramatica}

A Tabela~\ref{tab:gramatica} define a gramática. A hierarquia das regras é a forma da música: um \nt{Bloco} de 16 compassos tem duas
frases de 4 ciclos, e só o último ciclo de cada frase (\nt{Fim\_A}, \nt{Fim\_B}) pode ser
\tm{virada} ou \tm{breque}. Assim as marcas de forma caem sempre no fim dos compassos 8 e 16.
Nos demais ciclos, cada \nt{G} escolhe uma única coisa: a levada (\tm{base}), uma variante de
tamborim ou de pandeiro, ou uma frase de cuíca. Com isso, no máximo um instrumento varia por
ciclo. A Figura~\ref{fig:arvore} mostra a derivação de um bloco.

\begin{table}[t]
  \centering
  \footnotesize
  \renewcommand{\arraystretch}{0.9}
  \begin{tabular}{@{}r@{\;}c@{\;}l@{\qquad}r@{}}
    \toprule
    \nt{Música}   & $\to$  & \tm{cont} \nt{Blocos} \tm{final}  & $1$ \\
    \nt{Blocos}   & $\to$  & \nt{Bloco} \nt{Blocos}            & $p$ \\
                  & $\mid$ & \nt{Bloco}                        & $1-p$ \\
    \nt{Bloco}    & $\to$  & \nt{Frase\_A} \nt{Frase\_B}       & $1$ \\
    \nt{Frase\_A} & $\to$  & \nt{G} \nt{G} \nt{G} \nt{Fim\_A}  & $1$ \\
    \nt{Frase\_B} & $\to$  & \nt{G} \nt{G} \nt{G} \nt{Fim\_B}  & $1$ \\
    \nt{Fim\_A}   & $\to$  & \nt{G}                            & $1-v_a$ \\
                  & $\mid$ & \tm{virada}                       & $v_a$ \\
    \nt{Fim\_B}   & $\to$  & \tm{virada}                       & $1-b$ \\
                  & $\mid$ & \tm{breque}                       & $b$ \\
    \midrule
    \nt{G}        & $\to$  & \tm{base}                         & $1-\varepsilon-c$ \\
                  & $\mid$ & \nt{Var\_TB}                      & $\varepsilon/2$ \\
                  & $\mid$ & \nt{Var\_PD}                      & $\varepsilon/2$ \\
                  & $\mid$ & \nt{Var\_CU}                      & $c$ \\
    \nt{Var\_TB}  & $\to$  & \tm{tb1} $\mid$ \tm{tb2} $\mid$ \tm{tb3} & $0{,}40 \mid 0{,}35 \mid 0{,}25$ \\
    \nt{Var\_PD}  & $\to$  & \tm{pd1} $\mid$ \tm{pd2}          & $0{,}5 \mid 0{,}5$ \\
    \nt{Var\_CU}  & $\to$  & \tm{cu1}                          & $(1-v) + v/4$ \\
                  & $\mid$ & \tm{cu2} $\mid$ \tm{cu3} $\mid$ \tm{cu4} & $v/4$ cada \\
    \bottomrule
  \end{tabular}
  \caption{Regras de produção (axioma \nt{Música}): acima, a forma; abaixo, o conteúdo de
  cada ciclo. Na cuíca, com probabilidade $v$ sorteia-se uma das quatro frases e, senão,
  repete-se \tm{cu1}, que por isso soma $(1-v)+v/4$. \nt{Blocos} é truncado em
  $n_{\max}=10$.}
  \label{tab:gramatica}
\end{table}

\begin{figure}[t]
  \centering
  \begin{forest}
    for tree={font=\scriptsize, inner sep=1pt, s sep=7pt, l sep=5pt, l=0pt}
    [\nt{Bloco}
      [\nt{Frase\_A}
        [\nt{G} [\nt{Var\_PD} [\tm{pd1}, tier=t]]]
        [\nt{G} [\nt{Var\_TB} [\tm{tb1}, tier=t]]]
        [\nt{G} [\tm{base}, tier=t]]
        [\nt{Fim\_A} [\nt{G} [\nt{Var\_CU} [\tm{cu1}, tier=t]]]]
      ]
      [\nt{Frase\_B}
        [\nt{G} [\tm{base}, tier=t]]
        [\nt{G} [\tm{base}, tier=t]]
        [\nt{G} [\nt{Var\_CU} [\tm{cu1}, tier=t]]]
        [\nt{Fim\_B} [\tm{breque}, tier=t]]
      ]
    ]
  \end{forest}
  \caption{Derivação do 2º bloco da seed 7 (parâmetros padrão).}
  \label{fig:arvore}
\end{figure}

\subsection{Vocabulário: de terminal a arranjo}\label{sec:vocabulario}

Cada terminal vira o arranjo completo de um ciclo: \tm{base} é a levada da
Figura~\ref{fig:grade}; \tm{tb1--3} alteram 2 ou 3 casas do tamborim; \tm{pd1} desloca os
tapas e \tm{pd2} antecipa o grave do pandeiro; \tm{cu1--4} somam à base uma frase de cuíca;
\tm{breque} troca o 2º compasso por uma batida em uníssono e silêncio.

Como os padrões são pré-escritos, cada um foi avaliado frase a frase, tocado isolado e
comparado com sambas conhecidos. A cuíca deu mais trabalho, porque nos
sambas ela raramente segue um padrão fixo. Suas frases foram reescritas várias vezes, ouvindo
referências como ``Falador Passa Mal'' (Originais do Samba) e ``Taj Mahal'' (Jorge Ben Jor),
até ficar claro o seu papel: acentuar, com semicolcheias em espaços que não disputam com os
outros instrumentos, numa estrutura de pergunta e resposta em que o 2º compasso resolve a
frase do 1º, começando nos contratempos do tamborim. O pandeiro
também foi ajustado de ouvido, para encaixar as três pistas e a acentuação; por falta de mais
conhecimento do instrumento, suas variações só mudam a ordem de algumas batidas.

\subsection{Hiperparâmetros}\label{sec:hiper}

O principal hiperparâmetro é $\varepsilon$, que vai de uma levada fixa a uma percussão em
que os instrumentos se revezam em variações. Na cuíca, $c$
controla a presença e $v$ a variedade ($v=0$ faz de \tm{cu1} um motivo); $v_a$ e $b$
controlam as marcas de forma, e $p$, o tamanho: o número de blocos $n$ segue uma geométrica
truncada, com $\mathbb{E}[n] = (1-p^{\,n_{\max}})/(1-p) \approx 4{,}46$ para $p=0{,}8$, e
toda música tem $2+8n$ ciclos. Os padrões são $\varepsilon=0{,}15$, $c=0{,}25$, $v=0{,}5$,
$p=0{,}8$ e $v_a=b=0{,}3$, com $\varepsilon+c\le 1$; o andamento fica fora do método.

\subsection{Do símbolo ao áudio}\label{sec:audio}

Cada passo vira uma nota MIDI no instante $(16c+i)\Delta$, com
$\Delta = 60/(4\cdot\mathrm{BPM})$, mais um ruído gaussiano de 4\,ms, e a mesma seed
reproduz o mesmo \texttt{.mid}, gravado junto de um JSON com a derivação.
Com os instrumentos de percussão General MIDI, a música soou completamente
descaracterizada. A solução foi criar, com ajuda de IA, um sampler próprio que troca cada
nota MIDI por um recorte de instrumento gravado do BRID. Os ataques das faixas solo foram
recortados, separados por articulação e conferidos de ouvido, para garantir que cada
instrumento soasse limpo e audível. Ajustes de fade-in e fade-out foram decisivos na cuíca,
um instrumento de som contínuo que, com cortes secos, soava mecânico e pouco natural. Ainda
assim, a cuíca real não passa de forma binária do grave para o agudo, e o sampler, feito de
recortes de altura fixa, não reproduz essa transição.

% =====================================================================
\section{Resultados}\label{sec:resultados}

As três músicas entregues usam a seed 3, 90 BPM e os parâmetros padrão, salvo os
indicados\footnote{\url{https://github.com/Vipigal/gerador-backing-track-samba}, pasta
\texttt{output/}.}. A \texttt{musica\_1}
($\varepsilon=c=0$; 46\,s) é o controle: só sobram \tm{base}, \tm{virada} e \tm{breque}, e a
levada vira um loop em que apenas a forma varia. Na \texttt{musica\_2} ($\varepsilon=0{,}7$,
$c=0{,}3$, $v=0{,}3$; 1min07), tamborim e pandeiro variam em 54\% dos ciclos, com \tm{pd2}
quatro vezes seguidas, e a âncora quase some; a cuíca, com $v$ baixo, só repete \tm{cu1},
que soa como motivo. Na \texttt{musica\_3} ($\varepsilon=0{,}1$, $c=0{,}9$, $v=0{,}7$,
100 BPM; 3min15), nenhum ciclo é \tm{base}, e a cuíca, revezando as quatro frases em 71\%
dos ciclos, deixa de comentar a levada e vira uma camada constante. Entre esses extremos, o
padrão $\varepsilon=0{,}15$ mantém a âncora sem cair no
loop. Seeds diferentes geram músicas diferentes, o que os
testes verificam.

% =====================================================================
\section{Discussão e Trabalhos Futuros}\label{sec:discussao}

A linguagem gerada é \textbf{regular}: \nt{Blocos} só tem recursão à direita e o resto é
finito, então uma cadeia de Markov com estados de posição na frase geraria as mesmas
músicas. A gramática se justifica pela legibilidade: \nt{Fim\_A} e \nt{Fim\_B} \emph{são}
esses estados, com nome e significado musical.

Limitações: (a)~nada anuncia uma virada ou um breque, pois o fim da frase é sorteado sem
relação com o ciclo anterior; é um limite do desenho, não do formalismo, já que um terminal
de anúncio cabe na gramática, e uma cadeia de Markov não faria melhor;
(b)~a gramática não lembra escolhas anteriores, então o surdo é fixo e a cuíca pode repetir
a mesma frase em ciclos seguidos; (c)~a variação é combinatória sobre padrões escritos à mão;
(d)~falta à cuíca a transição entre grave e aguda.
Trabalhos futuros: anunciar as
marcas de forma, uma gramática de atributos (surdo e motivo da cuíca por música), recortes
de transição para a cuíca e microtiming medido no BRID.

% =====================================================================
% Fora do limite de páginas

\bibliography{referencias}

\end{document}
